\section*{\texorpdfstring{\S\CurrentExerciseGroup}{第{\CurrentExerciseGroup}节}}
\phantomsection
\addcontentsline{toc}{section}{第{\CurrentExerciseGroup}节习题}

\begin{enumerate}
\def\labelenumi{\arabic{enumi})}
\item
  设 E 是由函数 \(x \mapsto g(x) = \int_{0}^{x} f(t) dt\) （定义在 R 的区间 \([0,1]\) 上）所组成的向量空间，其中 f 取遍 \([0,1]\) 上的受控函数。设 P 是属于 E 的递增函数所成的集。试证 P 是一个凸锥，满足 E = P - P，\(P \cap (-P) = \{0\}\) ，并且 E 赋以由关系 \(g - h \in P\) 所定义的序结构后是一个 Riesz 空间。
\item
  设 \(I \subset \mathbf{R}\) 是一个紧区间。试证 \(\mathbf{R}^I\) 中由多项式函数（实系数）在 \(I\) 上的限制所组成的子空间是一个有向序空间，但不是 Riesz 空间。
\item
  \begin{enumerate}
  \def\labelenumii{\alph{enumii})}
  \tightlist
  \item
    设 E 是一个 Riesz 空间，H 是 E 的一个线性子空间。若对 H 中任意一对元素 x, y，\(\sup(x,y)\) 都属于 H，则称 H 为一个余格子子空间。此事成立当且仅当关系 \(x \in H\) 蕴涵 \(|x| \in H\) 。
  \end{enumerate}
\end{enumerate}

\begin{enumerate}
\def\labelenumi{\alph{enumi})}
\setcounter{enumi}{1}
\tightlist
\item
  设 I 是 R 中的一个紧区间，H 是 \(R^{I}\) 中由一次多项式 \(t \mapsto \alpha t + \beta\) 在 I 上的限制所组成的子空间。试证 H 不是 \(R^{I}\) 的余格子子空间，但 H 是一个 Riesz 空间（关于由 \(R^{I}\) 的结构诱导的结构）。
\end{enumerate}

\begin{enumerate}
\def\labelenumi{\arabic{enumi})}
\setcounter{enumi}{3}
\item
  设 E 是一个 Riesz 空间，H 是 E 的一个线性子空间。若关系 \(x \in H\) 与 \(|y| \leqslant |x|\) 蕴涵 \(y \in H\) ，则称 H 是 E 的一个孤立子空间。在此情形，设 P 是 E/H 中这样的元素 \(\dot{x}\) 所成的集：至少存在一个元素 \(x \geqslant 0\) 属于类 \(\dot{x}\) 。试证 P 是 E/H 上某个序结构的元素 \(\geqslant 0\) 全体所成的集，该序结构与 E/H 的向量空间结构相容，且对此结构 E/H 是一个 Riesz 空间（参见 A, VI, §1, 习题 4）。
\item
  \begin{enumerate}
  \def\labelenumii{\alph{enumii})}
  \tightlist
  \item
    若每个 \(x \in E\) ，只要使得 nx（n 为整数 \(\geqslant 0\) ）所成的集上有界，就必然 \(\leqslant 0\) ，则称序向量空间 E 是阿基米德的（A, VI, §1, 习题 31）。试证这一条件等价于：过 0 的任一平面与 E 的元素 \(\geqslant 0\) 所成的凸锥 P 的交是一个闭角域。试证为使一个序向量空间是阿基米德的，必须且只须它同构于某个完全格序空间的一个子空间（参见 A, VI, §1, 习题 31）。
  \end{enumerate}
\end{enumerate}

\begin{enumerate}
\def\labelenumi{\alph{enumi})}
\setcounter{enumi}{1}
\tightlist
\item
  设 F 是实值函数的 Riesz 空间 \(\mathbf{R}^{\mathbf{R}}\) （这些函数定义在 \(\mathbf{R}\) 上），并设 H 是有界函数所成的子空间 \(\mathcal{B}(\mathbf{R})\) （这些函数定义在 \(\mathbf{R}\) 上）；H 是 F 的一个孤立子空间（习题 4）。试证 Riesz 空间 \(\mathrm{E} = \mathrm{F} / \mathrm{H}\) （习题 4）不是阿基米德的：更确切地说，试证对 E 的每个元素 \(x \geqslant 0\) ，存在 \(y \in \mathrm{E}\) ，使得 \(y \geqslant nx\) 对每个整数 \(n \geqslant 0\) 成立。
\end{enumerate}

\begin{enumerate}
\def\labelenumi{\arabic{enumi})}
\setcounter{enumi}{5}
\item
  设 E 是一个完全格序空间，V 是 E 的一个线性子空间。试证：为使 E 中存在 V 的一个补 W 使 E 是 V 与 W 的序直和，必须且只须 V 是一个带；这时 W 必然等同于 E 中与 V 的每个元素都互奇异的元素所成的带。
\item
  设 E 是一个完全格序空间，H 是 E 的一个孤立子空间（习题 4）。若 \(B_{1}\) 与 \(B_{2}\) 是 E 中两个互余的带，试证 H 是 \(H_{1}=H\cap B_{1}\) 与 \(H_{2}=H\cap B_{2}\) 的序直和，且 Riesz 空间 E/H 是 Riesz 空间 \(B_{1}/H_{1}\) 与 \(B_{2}/H_{2}\) 的序直和。
\item
  设 \((\mathbf{E}_{\iota})\) 是一族完全格序空间，E 是诸 \(E_{\iota}\) 的乘积空间（它是完全格序的）。试证：若 B 是 E 中的一个带，则 B 的每个投影 \(\mathrm{B}_{\iota} = \mathrm{pr}_{\iota}(\mathrm{B})\) 是 \(E_{\iota}\) 中的一个带，且 B 等同于诸 \(B_{\iota}\) 的乘积。由此推出一个集合 A 到 R 中的映射所成的空间 \(R^{A}\) 中诸带的确定。试证在 A 上有界实值函数所成的空间 \(\mathcal{B}(A)\) 中，每个带都是某个带在 \(\mathcal{B}(A)\) 上的迹，而该带位于 \(R^{A}\) 中。
\item
  设 E 是一个完全格序空间。若 \(\mathfrak{F}\) 中存在一个上有界（相应地，下有界、有界）的集合，则称 E 上的滤子 \(\mathfrak{F}\) 是上有界的（相应地，下有界的、有界的）。对上有界（相应地，下有界）的滤子 \(\mathfrak{F}\) ，\(\mathfrak{F}\) 的上极限（相应地，下极限），记作 lim sup \(\mathfrak{F}\) （相应地，lim inf \(\mathfrak{F}\) ），定义为 E 的元素 \(\inf_{\mathrm{X}}(\sup \mathrm{X})\) （相应地，\(\sup_{\mathrm{X}}(\inf \mathrm{X})\) ），其中 X 取遍
\end{enumerate}

\(\mathfrak{F}\) 中上有界（相应地，下有界）的集合所成的集。若 \(\mathfrak{F}\) 有界，则 \(\liminf\mathfrak{F} \leqslant \limsup\mathfrak{F}\) ；称 \(\mathfrak{F}\) 关于 E 的序结构有一个极限或收敛，如果 \(\limsup\mathfrak{F} = \liminf\mathfrak{F}\) ；这时这两个元素的公共值记作 \(\lim\mathfrak{F}\) ，并称 \(\mathfrak{F}\) 收敛于 E 的这个元素。若有界滤子有一个极限，则每个更细的滤子有相同的极限（关于序结构）。

\begin{enumerate}
\def\labelenumi{\alph{enumi})}
\tightlist
\item
  为使有界滤子 \(\mathfrak{F}\) 关于序结构有一个极限，必须且只须 \(\inf_{\mathbf{X}} (\sup \mathbf{X} - \inf \mathbf{X}) = 0\) ，其中 \(\mathbf{X}\) 取遍 \(\mathfrak{F}\) 中有界的集合（可以利用下述事实：若在 E 中 \((x_{\alpha})\) 与 \((y_{\alpha})\) 是两个有下界的递减族，具有同一个右有向的指标集，则
\end{enumerate}

\[
\inf (x _ {\alpha} + y _ {\alpha}) = \inf x _ {\alpha} + \inf y _ {\alpha};
\]

为证明此公式，只需注意 \(\inf (x_{\alpha} + y_{\alpha})\leqslant x_{\beta} + y_{\gamma}\) 对每一对指标 \(\beta ,\gamma)\) 成立。

\begin{enumerate}
\def\labelenumi{\alph{enumi})}
\setcounter{enumi}{1}
\item
  设 A 是被滤子 G 过滤的集合，f 是 A 到 E 中的一个映射；若 \(f(\mathfrak{G})\) 是 E 中一个有极限（关于序结构）的有界滤子的基，则称 f 关于 G、按 E 的序结构有一个极限。试证：若 A 是一个有向序集，且 f 是 A 到 E 中的递增映射并在 A 上有上界，则 f 关于 A 的截面滤子有一个极限，等于 \(\sup f(x)\) 。
\item
  设 \((\mathbf{E}_{\iota})\) 是一族完全格序空间，E 是诸 \(E_{\iota}\) 的乘积完全格序空间。为使 E 上的有界滤子 F 关于序结构收敛，必须且只须对每个 \(\iota\) ，以 \(\Pr_{\iota}(\mathfrak{F})\) 为基的滤子在 \(E_{\iota}\) 中收敛；若 \(a_{\iota}\) 是其极限，则 \(a = (a_{\iota})\) 是 F 的极限。
\end{enumerate}

¶ 10) a) 设 E 是一个完全格序空间；E 上存在一个拓扑 \(\mathcal{T}_0(E)\) ，它是 E 的诸拓扑 \(\mathcal{T}\) 中最细的，使得每个按序结构意义（习题 9）收敛的有界滤子 \(\mathfrak{F}\) 对 \(\mathcal{T}\) 收敛于同一极限。

设 E 与 F 是两个完全格序空间，f 是 E 到 F 中的一个映射，使得对 E 上每个关于序结构收敛的有界滤子 F，\(f(\mathfrak{F})\) 是 F 上一个有界滤子的基，且该滤子关于序结构收敛于 \(f(\lim \mathfrak{F})\) 。试证在这些条件下，f 对拓扑 \(\mathcal{T}_{0}(\mathrm{E})\) 与 \(\mathcal{T}_{0}(\mathrm{F})\) 连续。

\begin{enumerate}
\def\labelenumi{\alph{enumi})}
\setcounter{enumi}{1}
\item
  由 a) 推出拓扑 \(\mathcal{T}_{0}(\mathrm{E})\) 与 E 的向量空间结构相容，且映射 \(x \mapsto |x|\) 对此拓扑连续。证明 \(\mathcal{T}_{0}(\mathrm{E})\) 与 E 的序向量空间结构相容，并由此推出此拓扑是 Hausdorff 的。
\item
  给定任一无穷集 A，设 \(\mathrm{E} = \mathcal{B}(\mathrm{A})\) 是 A 上有界实值函数所成的完全格序空间。设 \(\Phi\) 是定义在 A 上的数值函数 \(\varphi\) （有限与否皆可）的全体，这些函数满足 \(\varphi(t) > 0\) （对一切 \(t \in A\) ），且对每个整数 n \textgreater{} 0，\(t \in A\) 中使 \(\varphi(t) \leqslant n\) 者所成的集是有限的。对每个函数 \(\varphi \in \Phi\) ，设 \(V_{\varphi}\) 为 \(x \in E\) 中满足 \(|x| \leqslant \varphi\) 者所成的集；试证诸集 \(V_{\varphi}\) 构成一个拓扑 \(\mathcal{T}_{1}(\mathrm{E})\) 的 0 的基本邻域系，该拓扑与 E 的序向量空间结构相容，且对此拓扑 E 是 Hausdorff 的和完备的。试证拓扑 \(\mathcal{T}_{0}(\mathrm{E})\) 细于 \(\mathcal{T}_{1}(\mathrm{E})\) ，但严格粗于 A 上的一致收敛拓扑（由范数 \(\|x\| = \sup_{t \in A} |x(t)|\) 定义）（考虑元素 \(1 - \varphi_{X}\) ，它属于 \(\mathcal{B}(\mathrm{A})\) ，
\end{enumerate}

其中 X 取遍 A 的有限子集所成的集，\(\varphi_{X}\) 是 X 的特征函数）。由此推出，为使 E 的一个子集（关于序结构）有界，必须且只须它对拓扑 \(\mathcal{T}_{0}(E)\) 有界 (TVS, III, §1, No. 2)。试证 E 上每个有界且对拓扑 \(\mathcal{T}_{0}(E)\) 收敛的滤子关于序结构收敛于同一极限。最后，试证 E 上存在对 \(\mathcal{T}_{0}(E)\) 收敛但不有界的滤子。

¶ 11) 设 E 是一个完全格序空间，\((x_{\iota})_{\iota \in I}\) 是 E 的元素的一个族。对 I 的每个有限子集 H，令 \(s_{H} = \sum_{\iota \in H} x_{\iota}\) ；族 \((x_{\iota})_{\iota \in I}\) 称为关于 E 的序结构

可和的，如果映射 \(\mathrm{H} \mapsto s_{\mathrm{H}}\) 关于这一序结构、关于 I 的有限子集所成的有向序集 \(\mathfrak{F}(\mathrm{I})\) 有一个极限；这时这个极限 \(s\) 称为族 \((x_{\iota})_{\iota \in \mathrm{I}}\) 的和，记作 \(\sum_{\iota \in \mathrm{I}} x_{\iota}\) 。

\begin{enumerate}
\def\labelenumi{\alph{enumi})}
\item
  为使族 \((x_{\iota})_{\iota \in I}\) 可和，必须且只须：\(1^{\circ}\) 对 I 的每个有限子集 H，当 K 取遍 I 中与 H 不相交的有限子集所成的集时，诸 \(|s_{K}|\) 所成的集在 E 中有一个上确界 \(r_{\mathrm{H}}\) ；\(2^{\circ} \inf_{\mathrm{H}} r_{\mathrm{H}} = 0\) 当 H 取遍 \(\mathfrak{F}(I)\) （利用习题 9 a））。
\item
  把交换拓扑群中可和族的诸性质（GT, III, §5, 命题 2、3 与 Th. 2）推广到关于 E 的序结构的可和族。
\item
  设 \((x_{\iota})_{\iota \in I}\) 是 E 的元素 \(\geqslant 0\) 的一个族；为使该族可和，必须且只须有限部分和 \(s_{\mathrm{H}}\) 所成的集上有界（利用习题 9 b））。若 \((x_{\iota})_{\iota \in I}\) 可和，且 \((y_{\iota})_{\iota \in I}\) 是满足 \(0 \leqslant y_{\iota} \leqslant x_{\iota}\) （对一切 \(\iota\) ）的元素的一个族，试证族 \((y_{\iota})\) 可和，且 \(\sum_{\iota \in I} y_{\iota} \leqslant \sum_{\iota \in I} x_{\iota}\) ，其中等号仅当 \(x_{\iota} = y_{\iota}\) 对一切 \(\iota \in I\) 成立时成立。
\item
  设 \((x_{\iota})\) 是 E 的元素的一个族；试证若族 \((|x_{\iota}|)\) 关于序结构可和，则族 \((x_{\iota})\) 也可和。
\end{enumerate}

\begin{enumerate}
\def\labelenumi{\arabic{enumi})}
\setcounter{enumi}{11}
\tightlist
\item
  设 \(\mathbf{E}\) 是一个完全格序空间。试证存在一个族 \((u_{\iota})\) ，其元素 \(> 0\) 属于 \(\mathbf{E}\) ，使得 \(u_{\iota}\) 与 \(u_{\kappa}\) 对每一对不同的指标 \(\iota, \kappa\) 都互奇异，且对每个 \(x > 0\) （在 \(\mathbf{E}\) 中）至少存在一个指标 \(\iota\) 使 \(\inf(x, u_{\iota}) > 0\) （利用 Zorn 引理）。由此推出，对每个 \(x \geqslant 0\) ，存在唯一的一个族 \((x_{\iota})\) ，由元素 \(\geqslant 0\) （属于 \(\mathbf{E}\) ）组成，使得对每个 \(\iota, x_{\iota}\) 属于带 \(B_{\iota}\) （由 \(u_{\iota}\) 生成），且 \(x = \sum_{\iota} x_{\iota}\) （按 \(\mathbf{E}\) 的序结构）。
\end{enumerate}

（取 \(x_{\iota}\) 为 \(x\) 在带 \(B_{\iota}\) 中的分量）。反之，每个上有界的族 \((x_{\iota})\) ，由 E 的元素 \(\geqslant 0\) 组成，若满足 \(x_{\iota} \in B_{\iota}\) （对一切 \(\iota\) ），则在 E 中可和（习题 11）。

¶ 13 a) 设 E 是一个完全格序空间。若 \(u \neq 0\) 是 E 的任意一个元素，试证由这样的 \(x \in E\) 所成的集：存在一个整数 \(n\) 使得 \(|x| \leqslant n|u|\) ；该集是 E 的一个线性子空间 \(C_u\) ，它是 E 的一个完全格序的且余格子的子空间（习题 3）。对每个 \(x \in C_u\) ，用 \(\|x\|\) 表示标量 \(\lambda > 0\) （满足 \(|x| \leqslant \lambda |u|\) ）的下确界；试证 \(\|x\|\) 是 \(C_u\) 上的一个范数（利用 E 是阿基米德的这一事实（习题 5）），赋以此范数的 \(C_u\) 是完备的，且 \(\|x\| = \sup(\|x^+\|, \|x^-\|)\) 。

\begin{enumerate}
\def\labelenumi{\alph{enumi})}
\setcounter{enumi}{1}
\item
  设 \(\mathcal{I}\) 是 \(u\) 在空间 \(\mathrm{C}_u\) 的诸带中的分量所成的集；试证 \(\mathcal{I}\) 是由元素 \(c\) （属于 \(\mathrm{C}_u\) ，满足 \(0 \leqslant c \leqslant u\) 且 \(\inf(c, u - c) = 0\) ）组成的集（利用 Riesz 定理）。试证 \(\mathcal{I}\) 在赋范空间 \(\mathrm{C}_u\) 中是闭的，且对于由 \(\mathrm{C}_u\) 的序关系诱导的序关系，\(\mathcal{I}\) 是一个完备布尔代数（S, III, §1, 习题 11 与 17；只需证明若 \((c_\iota)\) 是 \(\mathcal{I}\) 的元素的一个族，则 \(\sup_{\iota} c_\iota\) 属于 \(\mathcal{I}\) ）。
\item
  设 \(x\) 是 \(\mathbf{C}_u\) 的任意一个元素；对每个 \(\lambda \in \mathbf{R}\) ，设 \(c(\lambda)\) 是 \(u\) 在 \(\mathbf{C}_u\) 的由 \((\lambda u - x)^+\) 生成的带中的分量。试证若 \(\lambda \leqslant \mu\) ，则 \(c(\lambda) \leqslant c(\mu)\) ；\(c(\lambda) = 0\) 对 \(\lambda < -\|x\|\) ，而 \(c(\lambda) = u\) 对 \(\lambda > \|x\|\) 。试证若 \(c \in \mathcal{I}\) 满足 \(c \leqslant c(\lambda)\) ，则 \(x\) 在由 \(c\) 生成的带中的分量 \(\leqslant \lambda c\) （注意由 \(c(\lambda)\) 生成的带与由 \((\lambda u - x)^+\) 生成的带是相同的，且 \(\lambda u - x\) 在此带中的分量等于 \((\lambda u - x)^+\) 在此带中的分量；由此推出 \(x\) 在这同一个带中的分量 \(\leqslant \lambda c(\lambda)\) ）。类似地，试证若 \(c \in \mathcal{I}\) 满足 \(c \leqslant u - c(\lambda)\) ，则 \(x\) 在由 \(c\) 生成的带中的分量 \(\geqslant \lambda c\) 。
\item
  已知（GT, II, §4, 习题 12）存在一个序结构同构 \(c \mapsto \theta_c\) ，它把布尔代数 \(\mathcal{I}\) 映到由一个紧全不连通空间 S 中既开又闭的集合的特征函数所组成的布尔代数上。试证关系 \(\sum_{i} \lambda_i c_i = 0\) 蕴涵函数 \(\sum_{i} \lambda_i \theta_{c_i}\) 在 S 上为 0（利用分解引理）；由这一注记与 c) 推出，映射 \(c \mapsto \theta_{c}\) 可以延拓为一个同构 \(x \mapsto \theta_{x}\) ，它把赋范空间 \(C_{u}\) 映到 S 上连续实值函数所成的赋范空间 \(\mathcal{C}(\mathrm{S};\mathbf{R})\) 上，使得 \((\theta_{x})^{+} = \theta_{x^{+}}\) 。（借助 c)，试证对每个 \(x \in C_{u}\) 与每个 \(\varepsilon > 0\) ，存在实数的一个递增序列 \((\lambda_{i})_{0 \leqslant i \leqslant n}\) ，使得 \(\lambda_{i} - \lambda_{i-1} \leqslant \varepsilon\)
\end{enumerate}

且 \(0 \leqslant x - \sum_{i=1}^{n} \lambda_{i-1} (c(\lambda_{i-1}) - c(\lambda_i)) \leqslant \varepsilon u\) 。

\begin{enumerate}
\def\labelenumi{\alph{enumi})}
\setcounter{enumi}{4}
\item
  为使 S 有限，必须且只须 \(C_{u}\) 是有限维的；由此（借助习题 12）推出每个有限维完全格序空间都同构于一个乘积空间 \(R^{n}\) （参见 §2, 习题 7）。
\item
  试证 S 是一个极不连通空间 (GT, I, §11, 习题 21)（利用 \(\mathcal{I}\) 是完备布尔代数这一事实）。极不连通紧空间称为 Stone 空间。试证若 X 是一个 Stone 空间，则对 X 上每个下半连续数值函数 \(f\) ，上半连续正则化 \(g\) （由 \(f\) 经正则化而得，GT, IV, §6, No. 2）是连续的（对每个 \(a < g(x)\) ，试证 \(x\) 不可能属于这样的 \(y\) 所成的（开）集的闭包：\(g(y) < a\) ；办法是注意 \(f(z) \leqslant a\) 在此集闭包的一点 \(z\) 处成立）。由此推出 \(\mathcal{C}(\mathrm{X};\mathbf{R})\) 这时是一个完全格序空间。试证当 Stone 空间 X 无穷时，一个上有界的集在 \(\mathcal{C}(\mathrm{X};\mathbf{R})\) 中的上确界未必等于它的上包络（考虑 X 中一个非闭的开集）；不过这两个函数在一个贫集的余集上相等（考虑使它们的差 \(\geqslant 1/n\) 的点所成的集）。
\item
  设 X 是一个 Stone 空间；试证若 f 是一个有界实值函数，定义在 X 的一个无处稠密子集 M 的余集上且在 X - M 上连续，则 \(f\) 可以连续地延拓到整个 X（利用 \(\mathcal{C}(X; \mathbf{R})\) 是完全格序的这一事实）。由此推出，若记 \(\mathcal{Q}(X; \overline{\mathbf{R}})\) 为 X 上（有限与否皆可的）连续数值函数 \(f\) 中满足 \(f(+\infty)\) 与 \(f(-\infty)\) 在 X 中无处稠密者所成的集，则它是一个完全格序向量空间。
\item
  试证 E 中由 u 生成的带 \(B_{u}\) （作为序空间）同构于 \(\mathcal{Q}(\mathrm{X};\overline{\mathbf{R}})\) 的一个孤立子空间（习题 4）（把元素 \(f \geqslant 0\) ——它属于 \(\mathcal{Q}(\mathrm{X};\overline{\mathbf{R}})\) ——看作诸元素 \(\inf(f,n)\) 的上确界）。
\end{enumerate}

\begin{enumerate}
\def\labelenumi{\arabic{enumi})}
\setcounter{enumi}{13}
\tightlist
\item
  设 E 是一个 Riesz 空间，带有一个与 E 的序向量空间结构相容的 Hausdorff 拓扑。
\end{enumerate}

\begin{enumerate}
\def\labelenumi{\alph{enumi})}
\item
  设 K 是 E 的一个紧子集，H 是 K 的一个对关系 \(\leqslant\) 有向的子集；试证 H 的截面滤子 F 收敛。（先证属于 K 的 H 的上界所成的集 L 非空且紧，再证 F 的接触点所成的集含于 L，最后证不可能存在两个不同的接触点。）
\item
  由 \(a\) ) 推出：若在 \(\mathbf{E}\) 中每个区间 \([a,b]\) 都是紧的，则 \(\mathbf{E}\) 是完全格序的。
\end{enumerate}

\setcounter{exercisegroup}{2}
\edef\CurrentExerciseGroup{\arabic{exercisegroup}}
